\documentclass[11pt]{article}
\usepackage[T1]{fontenc}
\usepackage[margin=1in]{geometry}
\usepackage{amsmath,amssymb,amsthm,booktabs,microtype}
\usepackage[hidelinks]{hyperref}
\newtheorem{theorem}{Theorem}[section]
\newtheorem{lemma}[theorem]{Lemma}
\newtheorem{proposition}[theorem]{Proposition}
\newtheorem{corollary}[theorem]{Corollary}
\theoremstyle{remark}
\newtheorem{remark}[theorem]{Remark}
\newcommand{\E}{\mathbb E}
\newcommand{\Prob}{\mathbb P}
\newcommand{\R}{\mathbb R}
\newcommand{\KL}{\operatorname{KL}}
\newcommand{\TV}{\operatorname{TV}}
\newcommand{\Lab}{\operatorname{Lab}}
\newcommand{\rank}{\operatorname{rank}}
\newcommand{\calC}{\mathcal C}
\numberwithin{equation}{section}
\title{Curvature from unlabelled two-dimensional causal orders}
\author{Alexander Kolpakov\\[0.5ex]
\small Center for Science, Technology, Engineering, and Mathematics\\
\small University of Austin, Austin, Texas, United States\\
\small \texttt{akolpakov@uaustin.org}\\[0.75ex]
\small The Wolfram Institute for Computational Foundations of Science\\
\small \texttt{akolpakov@wolframinstitute.org}}
\date{Research draft, 21 October 2026}
\begin{document}
\maketitle
\begin{abstract}
Consider independent events in a two-dimensional causal rectangle with a
positive smooth volume density, and retain only their unlabelled causal
order. We prove that every two-order representation recovers the empirical
null ranks to within $O(\log n/n)$, up to exchanging the two directions,
outside an event of probability $O(\log^2 n/n)$.
The proof combines classical modular decomposition with an empty-rectangle
argument for densities bounded above and below.
For a bounded $C^4$ class with fixed volume and uniform null marginals,
this gives an estimator of scalar curvature at the marginal centre with
mean squared error $O(n^{-2/5})$. A matching lower bound holds even
when all event coordinates are observed.
The permutation restrictions control the ambiguity of the latent
coordinates; the remaining two quotient frames form a genuine
two-element torsor. The statistical rate and the modular machinery
have established antecedents. The contribution addressed here is
their quantitative transfer to the partial-order observation model.
\end{abstract}

\section{The observation and the question}
Let $c$ be a positive density on $[0,1]^2$ with uniform marginals:
\begin{equation}\label{eq:copula}
 \int_0^1c(u,v)\,dv=\int_0^1c(u,v)\,du=1.
\end{equation}
The metric on the open square is
\begin{equation}\label{eq:metric}
 g_c=-2c(u,v)\,du\,dv.
\end{equation}
Its volume is one and its chronology is the strict product order.
Sample $Z_i=(U_i,V_i)$ independently with density $c$, then observe
only the unlabelled poset
\[
 i<_Pj\quad\Longleftrightarrow\quad U_i<U_j,\ V_i<V_j.
\]
A causal Hasse diagram supplies the same information after transitive
closure. Sampling labels, coordinates, an embedding and a null frame
are not observed.

Uniform marginals calibrate the null coordinates by marginal volume.
They do not assert that the two coordinates are independent.
For a physical rectangle of known volume $V$, the normalized model
has metric $-2Vc\,du\,dv$ and all curvatures below are divided by $V$.
The rectangle and its volume are part of the model.

The target point is $z_*=(1/2,1/2)$, invariant under exchange of the
two null directions. The scalar curvature, with the convention used
throughout these notes, is
\begin{equation}\label{eq:curvature}
 R_c(z_*)=\left.
 \left(\frac{2c_{uv}}{c^2}-\frac{2c_uc_v}{c^3}\right)\right|_{z_*}.
\end{equation}
Indeed $\Gamma^u_{uu}=\partial_u\log c$ and
$\Gamma^v_{vv}=\partial_v\log c$ are the only nonzero Christoffel
symbols; contracting the Ricci tensor gives
$R=2c^{-1}\partial_u\partial_v\log c$.

Order-based embedding and curvature recovery are established subjects.
Two-order models of causal rectangles already underlie
Brightwell--Henson--Surya~\cite{BHS}, including their discussion of
correlated null coordinates for a nonconstant conformal factor.
Henson~\cite{Henson} reconstructs approximate flat embeddings.
Benincasa--Dowker~\cite{BD} construct order-based curvature operators.
The copula--permutation correspondence is standard; see
Gr\"ubel~\cite{Grubel}. Tian~\cite{Tian} studies curvature from paired
null ranks, explicitly distinguishing those ranks from their partial-order
intersection. Our question is whether that intersection suffices at the
same statistical exponent on a full-support smooth class.

\section{The restrictions and their ambiguity}
For a finite set $V$, the bijections $\Lab(V)=\operatorname{Bij}(V,[n])$
form an $S_n$-torsor. Pairs of rank bijections form an
$S_n\times S_n$-torsor. The admissible subset for $P$ is
\[
 \mathcal R(P)=\{(r_1,r_2):
       x<_Py\ \Longleftrightarrow\
       r_1(x)<r_1(y),\ r_2(x)<r_2(y)\}.
\]
This subset is generally not itself a torsor.
Its members are the ordered two-realizers of $P$.
The history presentation in~\cite{History} also uses restrictions of
permutation torsors, there to describe token wirings and schedules.
Here the restrictions specify latent coordinate orders with a given
intersection; a causal sampling law is an additional assumption.

Fix one such pair. A common interval is a vertex set consecutive in both
orders. Let $b$ be the largest size of a proper common interval,
counting singletons. Relabel the first order increasingly to obtain
the usual permutation representation.

\begin{lemma}\label{lem:deterministic}
Suppose $b<n/2$. For any other ordered two-realizer
$(r'_1,r'_2)$, either directly or after exchanging its two orders,
\begin{equation}\label{eq:rankbound}
 \max_{x\in V}\max_{j=1,2}|r'_j(x)-r_j(x)|\le b-1.
\end{equation}
After contracting the maximal proper strong modules, exactly two
ordered quotient realizers remain. Exchange of the two orders acts
freely and transitively on them.
\end{lemma}
\begin{proof}
We use the classical modular representation of a permutation graph:
strong modules are precisely strong common intervals in every realizer,
and all realizers are obtained by composing realizers of the quotient
graphs in its modular decomposition. Prime quotients have four ordered
graph realizers, obtained by exchanging the orders and reversing both.
These facts, including the strong-module assertion, are stated in
Crespelle--Paul~\cite[Section 2.2, Proposition 1]{CP}.

Apply them to the incomparability graph of $P$.
The root of its modular decomposition is prime.
Otherwise its quotient is complete or empty, so the children occur
monotonically in the two orders. Each child has size at most $b$.
Remove an end child: the union of the remaining children is a proper
common interval with size at least $n-b>b$, a contradiction.

Let $B_1,\ldots,B_k$ be the root's children.
Each is consecutive in both orders of every realizer, and $|B_i|\le b$.
The prime quotient has four graph realizers.
Reversing both orders reverses the nonempty comparable relation of
the quotient poset. It therefore cannot preserve $P$.
The two remaining quotient realizers exchange the orders.
After making this global exchange if necessary, the blocks have exactly
the same order and the same rank ranges in both representations.
Only ranks within a block can change, by at most $|B_i|-1$.
This proves~\eqref{eq:rankbound} and the final assertion.
\end{proof}

The lemma gives a quantitative rank bound from classical realizer
decomposition~\cite{CP}. For uniqueness of prime two-realizers, see
Pouzet--Zaguia~\cite[Section 7.2]{PZ} and the finite theorem of
El-Zahar--Sauer cited there. The two ordered quotient frames form a
$\mathbb Z/2$-torsor; choices within the modules give the remaining
rank ambiguity.

\section{Localization for positive sampling densities}
Assume for this section that
\begin{equation}\label{eq:bounds}
 0<a\le c\le B<\infty
\end{equation}
and that~\eqref{eq:copula} holds. No differentiability is needed yet.

\begin{theorem}\label{thm:localize}
There are constants $K,C$ depending only on $a,B$ such that, for all
sufficiently large $n$, the following event has probability at least
$1-C\log^2n/n$:
every proper common interval of the sampled ranks has at most
$K\log n$ elements.
Consequently every ordered two-realizer of the observed poset agrees
with the empirical null ranks within $K\log n/n$ after normalization,
up to one global exchange of axes.
\end{theorem}
\begin{proof}
Set $\tau=A\log n/n$, with a sufficiently large constant $A$ depending
on $a$. We establish three events.

\emph{Empty rectangles.}
With failure probability $O(n^{-6})$, every open axis-aligned rectangle
of area greater than $\tau$ contains a sampled point.
To prove this, use the uniform grid with
$\lceil8/\tau\rceil$ subdivisions.
An empty rectangle of area greater than $\tau$ contains an empty
grid rectangle of area at least $\tau/2$: moving each of its four
sides inward by at most one mesh loses area at most $\tau/2$.
There are $O(\tau^{-4})$ grid rectangles.
Any specified one of area at least $\tau/2$ is empty with probability
at most $\exp(-an\tau/2)$.
A union bound proves the claim when $A$ is large enough.

\emph{Occupancy of short intervals.}
With failure probability $O(n^{-6})$, every interval of the $u$-axis
of length at most $4\tau$ contains at most $12n\tau$ sampled
$U_i$'s.
Such an interval is contained in one of the $O(\tau^{-1})$ intervals
$[j\tau,(j+6)\tau]\cap[0,1]$.
Uniform marginals bound the corresponding binomial mean by $6n\tau$.
The exponential Markov bound for a binomial variable gives
$\Prob(X\ge2\mu_0)\le\exp(-\mu_0/3)$ when its mean is at most $\mu_0$.
Take $\mu_0=6n\tau$ and use another union bound.

\emph{Corner exclusion.}
With failure probability at most $64Bn\tau^2$, no sample lies in
one of the four corner squares of side $4\tau$.
This follows from the density upper bound and a union bound over
the $n$ points.

On these three events, let $S$ be a proper common interval with at
least two elements. Write its coordinate bounding rectangle as
$I\times J$, with widths $s,t$.
There is no sample in
$I\times([0,1]\setminus J)$ or in
$([0,1]\setminus I)\times J$.
Each set is a union of two rectangles. The empty-rectangle event gives
\[
 s(1-t)\le2\tau,\qquad t(1-s)\le2\tau.
\]
If one width is at most $1/2$, the other is at most $4\tau$,
and substituting back gives both widths at most $4\tau$, for
$\tau\le1/8$.
If both widths exceed $1/2$, both are at least $1-4\tau$.
Any sample outside $S$ must then lie outside both $I$ and $J$,
and therefore in a corner square of side $4\tau$.
Corner exclusion contradicts properness of $S$.
Thus both widths are small. The occupancy event gives
$|S|\le12n\tau=12A\log n$.
Singletons satisfy the same bound.
The total failure probability is $O(\log^2n/n)$.
For large $n$, $12A\log n<n/2$; apply
Lemma~\ref{lem:deterministic}.
\end{proof}

For uniform random permutations, substantially sharper common-interval
results are classical: typically only size-two proper blocks remain.
See Bouvel--Chauve--Mishna--Rossin~\cite[Theorem 3]{BCMR}
and the earlier work cited there.
Theorem~\ref{thm:localize} uses a weaker logarithmic bound uniformly
over the positive copula class. It does not assume that the
sample permutation is uniform.

\section{An estimator from the order alone}
Fix $0<a<1<B$ and $M>0$. Let $\calC$ be the densities satisfying
\eqref{eq:copula}--\eqref{eq:bounds}, extending as $C^4$ functions
to the closed square, with
\[
 \max_{1\le|\nu|\le4}\|\partial^\nu c\|_\infty\le M.
\]
The model is genuinely nonparametric.
For example $1+\theta\sin(2\pi u)\sin(2\pi v)$ belongs to a suitable
fixed class for every $|\theta|\le\theta_0<1$.

Uniformly label the observed abstract poset, using its labelling torsor,
and apply any fixed rule that selects an ordered two-realizer.
Such a rule can use transitive orientation of the incomparability
graph, a standard permutation-graph algorithm~\cite{CP}.
The ranks $(\widehat U_i,\widehat V_i)=(r_1(i)/n,r_2(i)/n)$ are the
only coordinates used by the estimator.
Uniform labelling makes its law independent of any presentation
of the unlabelled input. It is not an assertion that all two-realizers
are equiprobable under the unknown sampling law.

Choose the even, compactly supported kernel
\[
 k(t)=\frac{315}{256}(1-t^2)^4\mathbf1_{\{|t|\le1\}} .
\]
It is $C^3$, integrates to one and has zero first moment.
For $\nu=(r,s)\in\{(0,0),(1,0),(0,1),(1,1)\}$ put
\begin{equation}\label{eq:kernel}
 T_\nu=\frac1{nh^{2+|\nu|}}
       \sum_{i=1}^n
       k^{(r)}\left(\frac{1/2-\widehat U_i}{h}\right)
       k^{(s)}\left(\frac{1/2-\widehat V_i}{h}\right).
\end{equation}
Project $T_{00}$ onto $[a,B]$ and each other $T_\nu$ onto $[-M,M]$;
denote the results by $\widetilde T_\nu$.
Define
\begin{equation}\label{eq:estimator}
 \widehat R=\frac{2\widetilde T_{11}}{\widetilde T_{00}^2}
       -\frac{2\widetilde T_{10}\widetilde T_{01}}
                         {\widetilde T_{00}^3}.
\end{equation}
Exchange of axes leaves~\eqref{eq:estimator} unchanged.
It is bounded in absolute value by
$R_{\max}=2M/a^2+2M^2/a^3$.

\begin{theorem}\label{thm:upper}
With $h=n^{-1/10}$, there is a constant $C=C(a,B,M)$ such that
\[
 \sup_{c\in\calC}\E_c|\widehat R-R_c(z_*)|^2\le Cn^{-2/5}
\]
for sufficiently large $n$.
The bound holds for every realizer-selection rule described above.
\end{theorem}
\begin{proof}
Let $E_n$ be the event of Theorem~\ref{thm:localize} and
$\epsilon_n=K\log n/n$.
Since~\eqref{eq:estimator} is invariant under exchanging axes,
we may align the selected ranks with the true empirical ranks on $E_n$.
Writing $F_{U,n},F_{V,n}$ for their empirical distribution functions,
set
\[
 \Delta_n=\max(\|F_{U,n}-\operatorname{id}\|_\infty,
                    \|F_{V,n}-\operatorname{id}\|_\infty).
\]
The uniform marginals and the Dvoretzky--Kiefer--Wolfowitz inequality
\cite{Massart} give
\begin{equation}\label{eq:dkw}
 \Prob(\Delta_n>t)\le4e^{-2nt^2},\qquad
 \E\Delta_n^4\le C/n^2.
\end{equation}
Dependence of $U_i$ and $V_i$ within a point does not affect either
one-dimensional empirical distribution bound.
On $E_n$, the aligned estimated coordinates satisfy
\begin{equation}\label{eq:coordinate}
 \max_i\|(\widehat U_i,\widehat V_i)-(U_i,V_i)\|_\infty
       \le\Delta_n+\epsilon_n.
\end{equation}

Let $T_\nu^0$ be~\eqref{eq:kernel} evaluated at the true sample
coordinates. Integration by parts, a symmetric-kernel Taylor expansion,
and $|\nu|\le2$ give
\begin{equation}\label{eq:oracle}
 |\E T_\nu^0-\partial^\nu c(z_*)|\le Ch^2,\qquad
 \operatorname{Var}(T_\nu^0)\le C/(nh^{2+2|\nu|}).
\end{equation}
The first bound uses derivatives of $c$ up to order $|\nu|+2\le4$.
For the second, the kernel is bounded by $Ch^{-2-|\nu|}$
and supported on a square of probability at most $4Bh^2$.

For large $n$, $\epsilon_n\le h/2$.
Restrict first to $\Delta_n\le h/2$.
Let $N_h$ count the true points in the square of radius $2h$ about
$z_*$.
A summand can change in~\eqref{eq:kernel} only if its true point lies
in this square. The bounded derivative of the scaled kernel and
\eqref{eq:coordinate} imply
\begin{equation}\label{eq:perturb}
 |T_\nu-T_\nu^0|
   \le Ch^{-3-|\nu|}(\Delta_n+\epsilon_n)N_h/n
                  \quad\text{on }E_n\cap\{\Delta_n\le h/2\}.
\end{equation}
Here $N_h$ is binomial with success probability at most $16Bh^2$.
When $nh^2\ge1$, expanding its second and fourth moments yields
\[
 \E(N_h/n)^2\le Ch^4,\qquad \E(N_h/n)^4\le Ch^8.
\]
Cauchy--Schwarz and~\eqref{eq:dkw} give
\[
 \E\{\Delta_n^2(N_h/n)^2\}\le Ch^4/n .
\]
These bounds do not assume independence of $N_h$ and $\Delta_n$.
Squaring~\eqref{eq:perturb}, with $|\nu|\le2$, gives
\[
 \E\bigl[|T_\nu-T_\nu^0|^2;
              E_n,\Delta_n\le h/2\bigr]
       \le C h^{-6}(n^{-1}+\epsilon_n^2).
\]

The true $c$ and its derivatives belong to the projection intervals,
so projecting cannot increase their estimation errors.
The map in~\eqref{eq:estimator} is Lipschitz on the product of these
intervals. On the complementary event its squared error is at most
$4R_{\max}^2$.
Combining the preceding bounds gives
\[
 \E|\widehat R-R_c(z_*)|^2
 \le C\left\{h^4+\frac1{nh^6}
          +\frac{\log^2n}{n^2h^6}
          +\frac{\log^2n}{n}+e^{-nh^2/2}\right\}.
\]
For $h=n^{-1/10}$ the first two terms are $n^{-2/5}$;
all the others are smaller for sufficiently large $n$.
\end{proof}

\begin{corollary}\label{cor:choice}
Let $\widehat R,\widehat R'$ use the same kernel, bandwidth
$h=n^{-1/10}$ and projection intervals, but any two possibly
randomized rules for selecting the two-realizer of the same input.
Under any coupling of these selections,
\[
 \sup_{c\in\calC}\E_c|\widehat R-\widehat R'|^2
       \le C\log^2n/n .
\]
Thus the choice between admissible realizers contributes a smaller
order than the minimax estimation error.
\end{corollary}
\begin{proof}
On $E_n$, align each pair with the true empirical ranks by its
global axis exchange. The resulting pairs differ by at most
$2\epsilon_n$ pointwise. On $\{\Delta_n\le h/2\}$ the argument
of~\eqref{eq:perturb} therefore gives
\[
 |T_\nu-T'_\nu|\le
       C h^{-3-|\nu|}\epsilon_n N_h/n .
\]
The binomial second-moment bound gives squared expectation at most
$C\epsilon_n^2h^{-6}$. Projection is nonexpansive and the curvature
map is Lipschitz on the projection intervals.
On the complement both scalar estimators are bounded by $R_{\max}$.
Hence their squared difference has expectation at most
\[
 C\{\epsilon_n^2h^{-6}+\log^2n/n+e^{-nh^2/2}\}.
\]
Substitution proves the claim. Every step is uniform in the
selection rules and their auxiliary randomness.
\end{proof}

\section{A matching information bound}
\begin{theorem}\label{thm:lower}
There is $c_0>0$ such that
\[
 \inf_{\widehat S}\sup_{c\in\calC}
       \E_c|\widehat S-R_c(z_*)|^2\ge c_0n^{-2/5}
\]
for sufficiently large $n$, even if $\widehat S$ observes all point
coordinates. Thus the unlabelled-order experiment attains the
full-coordinate minimax exponent.
\end{theorem}
\begin{proof}
Let $q$ be a smooth odd function supported in $(-1,1)$ with
$q'(0)\ne0$. For a small fixed $\gamma>0$ set
\[
 c_h(u,v)=1+\gamma h^4
          q((u-1/2)/h)q((v-1/2)/h).
\]
The integral of $q$ is zero, so both marginals remain exactly one.
For $h<1/2$ and sufficiently small $\gamma$, $c_h$ and the flat
density belong to $\calC$. Derivatives of order $j\le4$ are
bounded by a constant times $\gamma h^{4-j}$.
At $z_*$, the density equals one, both first derivatives vanish,
and $(c_h)_{uv}=\gamma h^2q'(0)^2$.
Thus the curvature difference is $2\gamma h^2q'(0)^2$.

For one full coordinate observation,
\[
 \KL(c_h\,du\,dv\|du\,dv)
 =\int(c_h\log c_h-c_h+1)\,du\,dv\le C\gamma^2h^{10}.
\]
For $n$ independent observations, the bound is $Cn\gamma^2h^{10}$.
Take $h=n^{-1/10}$ and $\gamma$ small enough that Pinsker's
inequality bounds total variation by $1/2$.
The nearest-curvature test associated with any estimator has sum
of testing errors at least $1/2$.
If the curvature difference is $d_h$, a classification error forces
squared estimation error at least $d_h^2/4$.
The larger of the two risks is therefore at least $d_h^2/16$,
of order $n^{-2/5}$.
Forgetting coordinates or labels cannot improve this testing bound.
\end{proof}

The estimator attains the classical two-dimensional second-derivative
rate using the partial order alone.

\begin{corollary}\label{cor:holder}
Replace the $C^4$ assumption by bounded derivatives through order two,
with each second derivative H\"older continuous of exponent
$0<\alpha\le1$ and seminorm at most $M$.
For this class, curvature at $z_*$ can be estimated from the
unlabelled order with minimax mean squared error of order
\[
 n^{-2\alpha/(6+2\alpha)}.
\]
The corresponding expected absolute error has minimax order
$n^{-\alpha/(6+2\alpha)}$.
\end{corollary}
\begin{proof}
The localization theorem is unchanged.
In~\eqref{eq:oracle} the bias is bounded by $Ch^\alpha$:
for the mixed derivative use its H\"older continuity directly,
and for the lower derivatives use the same symmetric Taylor expansion.
The perturbation and variance terms are unchanged.
Choose $h=n^{-1/(6+2\alpha)}$.
The term $\log^2n/n$ remains smaller than the resulting squared risk.
For the lower bound replace $h^4$ in $c_h-1$ by $h^{2+\alpha}$.
This preserves the H\"older bound and both marginals.
The curvature separation is of order $h^\alpha$ and the $n$-sample
relative entropy is at most $Cn h^{6+2\alpha}$.
The same two-point test proves the squared-risk lower bound.
Cauchy--Schwarz gives the absolute-error upper bound; the
nearest-curvature test gives its matching lower bound.
\end{proof}

This is also the exponent reported by Tian~\cite{Tian} for the
paired-rank experiment. Corollary~\ref{cor:holder} observes only
the intersection order. This comparison concerns the estimation
rate in the stated positive class, not finite-sample equivalence
of the two observation models.

\section{Examples and the role of the torsor}
For the Farlie--Gumbel--Morgenstern density
\[
 c_\theta(u,v)=1+\theta(1-2u)(1-2v),\qquad |\theta|\le\theta_0<1,
\]
the central curvature is $8\theta$.
This classical dependence family~\cite{Farlie} has comparable-pair
probability $1/2+\theta/9$. Indeed its lower-left cumulative
distribution is $uv\{1+\theta(1-u)(1-v)\}$, and comparability has
probability twice the integral of this function against $c_\theta$.
Expanding the integral gives the stated value.
Uniformly label an unlabelled sample of $2m$ events and pair the labels.
The $m$ comparability indicators $Y_j$ are then independent Bernoulli
variables with this parameter. Consequently
\[
 \widehat R_{\rm pair}
     =72\left(m^{-1}\sum_{j=1}^mY_j-\tfrac12\right)
\]
is unbiased for central curvature, with mean squared error at most
$1296/m$. Theorem~\ref{thm:upper} also applies
to nonparametric perturbations and to asymmetric positive copulas.
There is no requirement that a finite list of pair frequencies
identify the density.

For an asymmetric worked example, put
\[
 f(x)=1-2x,\qquad p(x)=6x^2-6x+1,
\]
and take
\[
 c(u,v)=1+\theta\{f(u)f(v)
             +\lambda f(u)p(v)+\mu p(u)f(v)\}.
\]
Both $f$ and $p$ have zero integral and absolute value at most one.
Thus the marginals are uniform, and
$|\theta|(1+|\lambda|+|\mu|)<1$ ensures positivity.
At the marginal centre the density is one, its first derivatives
are $\theta\lambda,\theta\mu$, and its mixed derivative is $4\theta$.
Consequently
\[
 R_c(z_*)=8\theta-2\theta^2\lambda\mu .
\]
For $(\theta,\lambda,\mu)=(2/5,3/10,1/2)$ the density lies between
$7/25$ and $43/25$, and the curvature is $394/125$.
The order-only estimator applies without knowing these parameters
or assuming this polynomial family.
Exchanging the axes interchanges $\lambda,\mu$ and preserves the target.

The torsor has two precise uses. Uniform labelling gives an invariant
randomized observation procedure. Restricting the product of the two
labelling torsors defines the admissible coordinate pairs.
Their modular quotient leaves a genuine two-element torsor; curvature
descends under its axis exchange. The proof also controls the fine
ambiguity left inside modules. Calling the whole realizer fibre a
torsor would be false.

Positivity throughout the rectangle is substantive.
With disconnected or vanishing-density support, large common blocks
can persist, and coordinate choices inside those blocks need not
agree up to one global axis exchange.
The theorem does not cover that case, noisy causal comparisons,
higher order dimension, or arbitrary rewrite ensembles.

The contribution is the localization of all realizers and the resulting
transfer of optimal curvature estimation to the partial-order observation
model. This uses classical realizer structure~\cite{CP,BBFGP,BCMR} and
copula ranks~\cite{Grubel}. Related work of Ghosh~\cite{Ghosh} treats
sample complexity and identifiability for counting observables on
finite causal sets.
The reconstruction question also occurs in Wolfram's discussion of
causal coordinates~\cite[Section 5.14]{Wolfram}; our theorem treats
the independent sampling model specified in Section~1.

\subsection*{Reproduction}
The exact checker \path{check_rank_ambiguity.py} enumerates all
ordered two-realizers for the $1\,840$ base permutations of sizes four
through seven that satisfy the deterministic interval hypothesis.
All $7\,456$ realizers obey the bound in
Lemma~\ref{lem:deterministic}.
The separate checker \path{check_curvature.py} checks the kernel's normalization and
boundary jets, the asymmetric example's marginals, and its curvature
directly from the metric. These finite checks supplement the proofs;
they do not test an asymptotic rate by simulation.
Run \texttt{make causal-order-paper} and
\texttt{make causal-order-check} from the repository root.
The exact outputs are included beside the source.

\begin{thebibliography}{99}
\raggedright
\bibitem{CP}
C. Crespelle and C. Paul,
\emph{Fully dynamic algorithm for recognition and modular decomposition
of permutation graphs}, Algorithmica \textbf{58} (2010), 405--432.
\href{https://webusers.i3s.unice.fr/~ccrespelle/publications/J02_DynPermu_long.pdf}{Author's full version}.
\bibitem{BBFGP}
F. Bassino, M. Bouvel, V. F\'eray, L. Gerin and A. Pierrot,
\emph{Dense and nondense limits for uniform random intersection graphs},
\href{https://arxiv.org/abs/2402.06394}{arXiv:2402.06394}.
\bibitem{BCMR}
M. Bouvel, C. Chauve, M. Mishna and D. Rossin,
\emph{Average-case analysis of perfect sorting by reversals},
\href{https://www.cecm.sfu.ca/~cchauve/Publications/DMAA_CPM09.pdf}{Author's full version}, Theorem 3.
\bibitem{Grubel}
R. Gr\"ubel, \emph{Ranks, copulas, and permutons},
Metrika \textbf{87} (2024), 155--182,
\href{https://doi.org/10.1007/s00184-023-00908-2}{doi:10.1007/s00184-023-00908-2}.
\bibitem{Massart}
P. Massart, \emph{The tight constant in the Dvoretzky--Kiefer--Wolfowitz
inequality}, Annals of Probability \textbf{18} (1990), 1269--1283,
\href{https://doi.org/10.1214/aop/1176990746}{doi:10.1214/aop/1176990746}.
\bibitem{Henson}
J. Henson, \emph{Constructing an interval of Minkowski space from a
causal set}, Classical and Quantum Gravity \textbf{23} (2006), L29--L35,
\href{https://arxiv.org/abs/gr-qc/0601069}{arXiv:gr-qc/0601069}.
\bibitem{BHS}
G. Brightwell, J. Henson and S. Surya,
\emph{A 2D model of causal set quantum gravity: the emergence of the continuum},
Classical and Quantum Gravity \textbf{25} (2008), 105025,
\href{https://arxiv.org/abs/0706.0375}{arXiv:0706.0375}.
\bibitem{PZ}
M. Pouzet and I. Zaguia,
\emph{Hereditary classes of ordered sets of width at most two},
\href{https://arxiv.org/abs/2112.02633}{arXiv:2112.02633}, Section 7.2.
\bibitem{BD}
D. M. T. Benincasa and F. Dowker,
\emph{The scalar curvature of a causal set},
Physical Review Letters \textbf{104} (2010), 181301,
\href{https://arxiv.org/abs/1001.2725}{arXiv:1001.2725}.
\bibitem{Farlie}
D. J. G. Farlie, \emph{The performance of some correlation
coefficients for a general bivariate distribution},
Biometrika \textbf{47} (1960), 307--323,
\href{https://doi.org/10.1093/biomet/47.3-4.307}{doi:10.1093/biomet/47.3-4.307}.
\bibitem{History}
A. Kolpakov, \emph{Causal histories, observable states, and cubical
geometry}, accompanying research manuscript, 2026,
\href{https://github.com/WolframInstitute/graph_of_graphs/tree/main/history-state-paper}{repository source and PDF}.
\bibitem{Tian}
C. Tian, \emph{Recovering Lorentzian curvature from finite volumes
and paired ranks}, preprint posted 16 September 2026,
\href{https://ssrn.com/abstract=7472426}{SSRN 7472426}.
\bibitem{Ghosh}
A. Ghosh, \emph{Sample-complexity bounds and an exact bulk-sector
channel obstruction for curvature estimation on finite causal sets},
11 August 2026,
\href{https://zenodo.org/records/21931074}{Zenodo, version 4}.
\bibitem{Wolfram}
S. Wolfram, \emph{A Class of Models with the Potential to Represent
Fundamental Physics}, 2020,
\href{https://www.wolframphysics.org/technical-introduction/the-updating-process-for-string-substitution-systems/foliations-and-coordinates-on-causal-graphs/}{Section 5.14}.
\end{thebibliography}
\end{document}
